% ============================================================================
% Visão Geral do Projeto de TCC — Sistema IoT para Monitoramento da Presença
% de Estudantes em Transporte Escolar Rural por Reconhecimento de Voz e Vídeo
%
% Aluna: Fernanda Andrade Bicalho da Cruz
% Orientador: Janier Arias García — Coorientador: Frederico Gualberto Ferreira Coelho
% Entrega: 18/10/2026 (adiada de 11/10/2026) (Atividades 1.1 / 1.2 do TCC I)
% Versão 01.00 finalizada em 18/10/2026.
%
% Estrutura segue o modelo oficial do curso:
%   docs/modelos/modelo_visao_geral_tcc.doc
% (Introdução -> Visão Geral do Trabalho -> Metodologia e Cronograma ->
%  Referências bibliográficas). As instruções do modelo foram omitidas,
%  conforme o próprio modelo pede para a versão entregue.
%
% Para compilar: envie este arquivo, "referencias.bib" e a pasta "imagens/"
% juntos para um projeto no Overleaf (mesmo procedimento de monografia_tcc1.tex).
%
% Antes de entregar:
%   - conferir as referências marcadas com "% CONFERIR" em referencias.bib;
%   - os parágrafos de voz e de trabalhos correlatos estão sem \cite: as
%     referências candidatas estão nos comentários "% CITAR:" logo abaixo deles
%     (ver docs/lacunas_revisao_bibliografica.md) e entram depois de lidas.
% ============================================================================

\documentclass[
    12pt,
    oneside,
    a4paper,
    article,
    english,
    brazil
    ]{abntex2}

% ----------------------------------------------------------------------------
% Pacotes
% ----------------------------------------------------------------------------
\usepackage{lmodern}
\usepackage[T1]{fontenc}
\usepackage[utf8]{inputenc}
\usepackage{indentfirst}
\usepackage{xcolor}
\usepackage{graphicx}
\usepackage{microtype}
\usepackage{float}
\usepackage{tabularx}
\usepackage{booktabs}

\usepackage[alf]{abntex2cite}

% Marca trechos que ainda dependem de informação/decisão — remover antes da entrega.
\newcommand{\pendente}[1]{\textcolor{red}{\textbf{[PENDENTE: #1]}}}

% ----------------------------------------------------------------------------
% Dados do documento (equivalentes aos campos de autotexto do modelo .doc)
% ----------------------------------------------------------------------------
\newcommand{\nomeprojeto}{Sistema IoT para Monitoramento da Presença de Estudantes em Transporte Escolar Rural por Reconhecimento de Voz e Vídeo}
\newcommand{\nomealuna}{Fernanda Andrade Bicalho da Cruz}
\newcommand{\versaodoc}{01.00}
\newcommand{\datadoc}{18 de outubro de 2026}

\titulo{Visão Geral do Projeto -- \nomeprojeto}
\autor{\nomealuna}
\local{Belo Horizonte}
\data{2026}

\makeatletter
\hypersetup{
    pdftitle={\@title},
    pdfauthor={\@author},
    pdfcreator={LaTeX com abnTeX2},
    colorlinks=true,
    linkcolor=blue,
    citecolor=blue,
    urlcolor=blue,
    bookmarksdepth=4
}
\makeatother

% ============================================================================
\begin{document}

\selectlanguage{brazil}
\frenchspacing

% ----------------------------------------------------------------------------
% Capa (conforme modelo: UFMG / Curso / TCC I / título / versão / data / aluno)
% ----------------------------------------------------------------------------
\begin{titlepage}
\begin{center}
    \includegraphics[width=3.5cm]{imagens/ufmg_logo.png}

    \vspace{0.5cm}
    {\large Universidade Federal de Minas Gerais}\\
    {\large Curso de Engenharia de Sistemas}\\
    {\large Trabalho de Conclusão de Curso I}
\end{center}

\vspace{4cm}
\begin{flushright}
    {\LARGE\textbf{Visão Geral do Projeto}}\\[0.5cm]
    \begin{minipage}{11cm}
        \raggedleft\large \nomeprojeto
    \end{minipage}
\end{flushright}

\vfill
\begin{center}
    {\large Aluna: \nomealuna}\\[0.2cm]
    Orientador: Janier Arias García\\
    Coorientador: Frederico Gualberto Ferreira Coelho
\end{center}
\end{titlepage}

% ----------------------------------------------------------------------------
% Histórico de revisões
% ----------------------------------------------------------------------------
\begin{center}
    {\large\textbf{Histórico de revisões}}
\end{center}

\begin{table}[H]
\centering
\small
\begin{tabularx}{\textwidth}{|c|c|p{4.2cm}|X|}
    \hline
    \textbf{Versão} & \textbf{Data} & \textbf{Autor} & \textbf{Descrição} \\
    \hline
    01.00 & 12/out/2026 & \nomealuna & Versão inicial, entregue em 18/10/2026. \\
    \hline
     & & & \\
    \hline
\end{tabularx}
\end{table}

\newpage

% ----------------------------------------------------------------------------
% Sumário
% ----------------------------------------------------------------------------
\pdfbookmark[0]{\contentsname}{toc}
\tableofcontents*
\newpage

% ============================================================================
% 1. INTRODUÇÃO
% ============================================================================
\section{Introdução}

\subsection{Propósito}

Este documento especifica o trabalho a ser desenvolvido pela aluna de
Engenharia de Sistemas \nomealuna, fornecendo ao professor orientador e ao
professor da disciplina as informações necessárias para sua avaliação e
aprovação. Além disso, o documento servirá como base para o desenvolvimento do
projeto ao longo dos dois semestres de sua realização e também pode servir para
gerar partes dos textos das monografias de TCC~I e TCC~II.

\subsection{Público Alvo}

Este documento se destina à aluna e aos profissionais envolvidos na orientação
de seu trabalho de conclusão de curso.

\subsection{Motivação e justificativa}

\textbf{Problema.} O transporte escolar rural é essencial para garantir o
acesso à educação em regiões afastadas, mas o controle de quais alunos
embarcam e desembarcam dos ônibus ainda depende, em geral, de processos
manuais. Esse registro é sujeito a falhas, difícil de auditar e não fornece
dados confiáveis sobre a operação das rotas.

\textbf{Importância.} A falta de um registro confiável tem consequências em três
frentes. Na segurança dos estudantes, não é possível afirmar com certeza se um
aluno embarcou, se desembarcou no ponto correto ou se ficou para trás. Na
transparência da gestão pública, os dados são difíceis de serem coletados e
auditados. Na eficiência
operacional, a falta de dados impede que os gestores identifiquem rotas
subutilizadas, atrasos recorrentes ou o custo real por aluno transportado.

\textbf{Benefícios esperados.} O trabalho propõe um dispositivo de Internet das
Coisas (IoT) embarcado no ônibus que identifica automaticamente cada aluno no
embarque e no desembarque, por reconhecimento facial e de voz, sem exigir
nenhuma ação do motorista e com o mínimo de ação do aluno. Os eventos de presença são enviados ao
sistema TRANSCOLAR RURAL, já utilizado na gestão do transporte escolar. Com
isso, espera-se: (i) maior segurança física dos estudantes, com alertas
imediatos de falha de reconhecimento ou embarque irregular; (ii) registros
auditáveis de presença, que dão transparência ao uso dos recursos públicos; e
(iii) relatórios de presença, localização, tempo de viagem e eficiência das
rotas para os gestores escolares. O sistema é projetado desde o início com
foco em privacidade: nenhuma imagem ou áudio dos alunos é armazenado, apenas
o evento de presença resultante.

\subsection{Objetivos}

O \textbf{objetivo geral} deste trabalho é projetar e implementar um dispositivo
IoT embarcado em ônibus escolares rurais que registre a presença dos alunos
por reconhecimento de voz e vídeo, integrando os dados gerados ao sistema
TRANSCOLAR RURAL, de modo a aumentar a segurança dos estudantes e a eficiência
operacional do transporte escolar.

Os \textbf{objetivos específicos} são:

\begin{enumerate}
    \item desenvolver um protótipo de hardware IoT dotado de câmera e microfone;
    \item implementar algoritmos de reconhecimento facial e de voz adequados às
          condições reais de operação no interior de um ônibus;
    \item integrar o dispositivo ao sistema TRANSCOLAR RURAL por meio de rede móvel;
    \item viabilizar a geração de relatórios inteligentes voltados a gestores escolares;
    \item avaliar o desempenho, a segurança e a confiabilidade do sistema proposto.
\end{enumerate}

\subsection{Local de realização}

O trabalho é realizado no Departamento de Engenharia Eletrônica (DELT) da
Escola de Engenharia da Universidade Federal de Minas Gerais. O DELT atua no ensino e na
pesquisa em eletrônica, sistemas embarcados, processamento de sinais e
inteligência computacional, áreas em que este projeto se insere: ele combina
hardware embarcado, processamento de imagem e áudio e aprendizado de máquina
aplicado a um problema de gestão pública.

\subsection{Visão geral do documento}

\begin{itemize}
    \item A Seção~\ref{sec:visao-geral} apresenta uma visão geral do trabalho,
          caracterizando o seu escopo e efetuando uma breve revisão bibliográfica
          sobre o assunto, permitindo situar claramente o trabalho proposto em
          relação a outros apresentados na literatura;
    \item A Seção~\ref{sec:metodologia} descreve a metodologia a ser empregada
          no trabalho a ser desenvolvido;
    \item A seção \textit{Referências} apresenta as referências bibliográficas
          utilizadas.
\end{itemize}

% ============================================================================
% 2. VISÃO GERAL DO TRABALHO
% ============================================================================
\section{Visão Geral do Trabalho}
\label{sec:visao-geral}

O trabalho consiste em um dispositivo embarcado instalado na porta do ônibus
escolar, composto por câmera, microfone e uma unidade de processamento local.
Quando o aluno embarca ou desembarca, o dispositivo detecta o rosto, compara-o
com as referências biométricas cadastradas e, opcionalmente, usa a voz como
reforço. Todo o reconhecimento é feito no próprio dispositivo (computação de
borda), sem depender de processamento em nuvem. O resultado é indicado na
hora ao motorista por um LED e um sinal sonoro, e gera um evento de
presença (aluno, ônibus, rota, horário e tipo de evento). A imagem e o áudio
capturados são descartados.

Os eventos são enviados ao backend e ao TRANSCOLAR RURAL pela rede móvel, por
meio do protocolo MQTT. Como a conectividade nas rotas rurais é intermitente,
o dispositivo mantém uma fila local e reenvia os eventos quando a conexão volta,
sem perda nem duplicação. A partir desses dados, o sistema gera relatórios de
presença, localização, tempo de viagem, segurança, eficiência operacional e
qualidade do serviço para os gestores escolares.

O ciclo completo se divide em duas etapas. O \textbf{TCC~I} cobre a
especificação e a concepção: levantamento de requisitos, revisão bibliográfica
com avaliação empírica de algoritmos, projeto arquitetural e análise de
impactos. O \textbf{TCC~II} cobre a realização: implementação do software
embarcado, integração com o TRANSCOLAR RURAL, simulação, teste piloto em
ônibus real e validação das métricas.

\subsection{Escopo do trabalho}

\textbf{O que se pretende fazer.} Especificar, projetar, implementar e validar,
em um piloto com um ônibus, um sistema que contempla:

\begin{itemize}
    \item reconhecimento facial no embarque e no desembarque, com
          reconhecimento de voz como reforço;
    \item retorno local ao motorista por LED e sinal sonoro, sem tela nem login;
    \item registro de eventos de presença com funcionamento offline e
          sincronização posterior, sem perda nem duplicação;
    \item integração com o TRANSCOLAR RURAL via MQTT/API;
    \item relatórios para gestores escolares;
    \item cadastro de alunos, rotas, ônibus e referências biométricas, com
          registro e revogação do consentimento dos responsáveis.
\end{itemize}

As principais metas quantitativas definidas no levantamento de requisitos são:
acurácia de reconhecimento superior a 90\%, tempo de no máximo 2~segundos por
aluno entre a detecção e o retorno local, e capacidade para cerca de 20 alunos
por ônibus no piloto.

\textbf{O que não se pretende fazer.} Não fazem parte do escopo: interfaces
(aplicativo, tela ou login) para motoristas, monitores ou responsáveis; a coleta
do consentimento dos responsáveis, que é um processo administrativo da escola
(o sistema apenas o registra); o desenvolvimento ou a alteração do próprio
TRANSCOLAR RURAL; e o uso de dados reais de alunos durante a prototipagem, que
utiliza apenas bases públicas de adultos.

\textbf{Inserção em um projeto maior.} O trabalho complementa o TRANSCOLAR
RURAL, sistema já utilizado na gestão do transporte escolar rural. O dispositivo
proposto funciona como uma nova fonte automática de dados de presença para esse
sistema, que hoje depende de registro manual. Caso a interface com esse
sistema não esteja disponível para integração com o projeto em desenvolvimento
(esse ponto ainda será definido), a integração será validada contra um serviço
simulado que segue o mesmo formato de mensagem; como o envio fica isolado em um
único componente do software embarcado, a troca pelo serviço real não altera o
restante do sistema.

\textbf{Situação atual e situação a ser alcançada.} Atualmente, o controle de
presença nos ônibus é manual, sem registro auditável e sem dados operacionais
consolidados. Ao final do trabalho, espera-se um protótipo validado em um
ônibus real que registre automaticamente o embarque e o desembarque de cada
aluno, alerte em caso de falha ou embarque irregular, funcione mesmo sem sinal
de rede e entregue esses dados ao TRANSCOLAR RURAL. Tudo isso respeitando a
Lei Geral de Proteção de Dados Pessoais (LGPD) \cite{brasil2018lgpd}, por se
tratar
de dados biométricos de menores de idade.

\subsection{Revisão bibliográfica}

Os primeiros sistemas automáticos de reconhecimento facial tratavam a imagem
inteira como um vetor de pixels. No Eigenfaces, \citeonline{turk1991eigenfaces}
usam análise de componentes principais (PCA) para descrever cada rosto por
poucos coeficientes, e o próprio artigo mostra o limite da ideia: o acerto
ficou perto de 96\% quando só a iluminação variava, mas caiu para 64\% quando
variava o tamanho do rosto na imagem. O Fisherfaces
\cite{belhumeur1997eigenfaces} troca a PCA pela análise discriminante linear,
que usa os rótulos do treino para separar as pessoas, e não as condições de
luz. Já o LBPH \cite{ahonen2004face} descreve o rosto pela textura local de
cada região, com histogramas de padrões binários. São métodos baratos, que
rodariam com folga em qualquer placa embarcada, mas a variação de luz e de
ângulo que os prejudica é exatamente o que se espera na porta de um ônibus.
Por isso aparecem aqui só como ponto de partida histórico.

A mudança veio com as redes profundas treinadas para produzir um vetor de
características (\textit{embedding}) em que a distância entre dois vetores já
mede a semelhança entre dois rostos. O FaceNet \cite{schroff2015facenet}
aprende esse espaço com a \textit{triplet loss}, que compara trios de imagens
(duas da mesma pessoa e uma de outra), e chegou a 99,63\% de acurácia na LFW.
Para este projeto, a consequência prática é que cadastrar um aluno novo exige
só guardar o \textit{embedding} da sua foto, sem treinar a rede de novo. Os
trabalhos seguintes passaram a medir a semelhança pelo ângulo entre os
vetores. O ArcFace \cite{deng2019arcface} soma uma margem fixa ao ângulo entre
o \textit{embedding} e o centro da classe correta e chegou a 99,83\% na LFW. O
SFace \cite{zhong2021sface} suaviza essa otimização com curvas sigmoides, para
que fotos ruins e rótulos errados da base de treino não dominem o
aprendizado. No benchmark deste trabalho (Tabela~\ref{tab:benchmark}), o
ArcFace também foi o modelo mais preciso, enquanto o SFace, com apenas 5~MB,
foi de longe o mais rápido.

Antes de qualquer comparação, o rosto precisa ser encontrado e alinhado, e
essa etapa acabou pesando tanto quanto o modelo. O RetinaFace
\cite{deng2020retinaface} prevê, numa só passada da rede, a caixa do rosto e
cinco pontos de referência, e lida bem com rostos pequenos. Com
ele, a acurácia do ArcFace no benchmark subiu de 92,97\% para 96,61\%, ao custo
de cerca de 0,6~s por imagem na CPU. O YuNet \cite{wu2023yunet} vai no sentido
oposto: tem cerca de 76 mil parâmetros, código aberto e roda em poucos
milissegundos na CPU, com precisão menor que a do RetinaFace nos rostos mais
difíceis. Se essa perda é aceitável dentro do ônibus é uma das perguntas que
a próxima rodada do benchmark deve responder.

Para comparar os modelos em condições iguais, foi usada a biblioteca DeepFace
\cite{serengil2020lightface}, que aplica o mesmo pipeline de detecção,
alinhamento e verificação a todos eles. Os autores da biblioteca lembram que
o alinhamento muda o resultado: no FaceNet, ele elevou a acurácia na
LFW de 98,87\% para 99,63\%. Os testes usam apenas uma base pública de
adultos, a \textit{Labeled Faces in the Wild} (LFW) \cite{huang2007lfw}, com
13.233 fotos de 5.749 pessoas tiradas de notícias, em condições não
controladas.

Para levar os eventos ao servidor, o projeto usa o MQTT \cite{oasis2019mqtt},
um protocolo leve de publicação e assinatura pensado para redes instáveis. No
nível de qualidade de serviço 1, o \textit{broker} confirma o recebimento de
cada mensagem, e é essa confirmação que permite ao dispositivo apagar um
evento da fila local só depois de entregue. Numa rota rural, em que o sinal
de celular cai com frequência, isso faz diferença.

A voz entra como reforço do rosto. No TCC~I ela fica restrita à revisão
bibliográfica, e a escolha e a avaliação do modelo de reconhecimento de
locutor ficam para o TCC~II. A lógica é a mesma do reconhecimento facial. No
ECAPA-TDNN \cite{desplanques2020ecapa}, que superou os sistemas baseados em
x-vector e em ResNet da época, uma rede extrai um \textit{embedding} de cada
trecho de fala, e dois trechos são comparados pela distância de cosseno entre
os vetores, como se faz com os rostos. A versão menor do modelo tem
6,2 milhões de parâmetros e taxa de erro igual (EER) de 1,01\% no VoxCeleb1.
Essa base \cite{nagrani2017voxceleb}, com mais de 100 mil falas de 1.251
celebridades tiradas de vídeos do YouTube, tem para a voz o papel que a LFW
tem para o rosto: as gravações vêm de ambientes reais, como entrevistas,
estádios e tapetes vermelhos, com ruído e várias pessoas falando. O ruído de
um ônibus, porém, é outro: o motor, a estrada e a conversa dos próprios
alunos devem piorar a taxa de acerto. O ECAPA-TDNN já é treinado com ruído,
fala de fundo e reverberação artificiais, mas isso não garante o desempenho
dentro do ônibus, e o ruído vai pesar tanto na escolha do modelo quanto na do
microfone direcional.
% CITAR: versão leve para borda (Chen et al., APSIPA ASC 2023), para sustentar
% que a voz caberia no mesmo hardware do rosto; a frase sobre modelos compactos
% em borda e microcontroladores saiu até essa referência entrar.
% Ver docs/lacunas_revisao_bibliografica.md, seção 1.2.

Há trabalhos anteriores com a mesma motivação. Em ônibus escolares,
\citeonline{garg2021onibus} registram o embarque e o desembarque com cartões
RFID e avisam os pais por SMS ou e-mail. Como o cartão exige uma ação do aluno
e pode ser usado por outra pessoa, os autores acrescentam uma câmera ESP32-CAM
com reconhecimento facial para confirmar quem embarca. Em sala de aula, o trabalho mais próximo deste é o de
\citeonline{duarte2025classic}, que usa a mesma combinação de Raspberry Pi,
DeepFace e ArcFace. Os números deles ajudam a dimensionar o problema: o
reconhecimento levou cerca de 4,5~s no próprio Raspberry Pi~4 e caiu para
cerca de 2~s quando foi transferido a um computador na rede local. O acerto
foi de 85,7\% entre 14 voluntários, abaixo dos 90\% que o mesmo par de detector
e modelo havia obtido com fotos.

O que este trabalho acrescenta é juntar essas peças sob as restrições do
transporte rural. O reconhecimento roda inteiro no dispositivo, sem depender
de rede, ao contrário do CLASSIC, que precisa de um computador na mesma rede
Wi-Fi para ser rápido. Quando o sinal cai, os eventos esperam numa fila local
e são enviados depois, sem perda nem duplicação. Nenhuma imagem é guardada,
enquanto o CLASSIC armazena todas as fotos capturadas, um ponto sensível pela
LGPD quando os dados são de crianças. E os eventos alimentam um sistema
público que já existe, o TRANSCOLAR RURAL.
% Frases sem fonte na pasta foram retiradas em 22/10/2026 (Raspberry Pi em
% ônibus e "menos de 1 s"). Candidatas para a monografia: ICCMC 2023, IEEE 2019
% e Surantha e Sugijakko (2024); ver docs/lacunas_revisao_bibliografica.md, seção 1.1.

% ============================================================================
% 3. METODOLOGIA E CRONOGRAMA
% ============================================================================
\section{Metodologia e Cronograma}
\label{sec:metodologia}

\subsection{Atividades do TCC I}

\textbf{Atividade 1.1 -- Levantamento de requisitos.} Os requisitos funcionais
e não funcionais foram levantados a partir da proposta do projeto e de
decisões tomadas com a orientação. Eles estão consolidados em um documento de
requisitos com identificadores rastreáveis até os objetivos específicos. Entre
as decisões já tomadas estão: não armazenar imagem nem áudio dos alunos; tempo
máximo de 2~s por reconhecimento; retorno local por LED e bipe; retenção dos
eventos no dispositivo apenas até a sincronização e, no servidor, pelo ano
letivo mais um ano; registro do consentimento dos responsáveis como
pré-requisito para o cadastro biométrico; e, entre os dois tipos de erro de
reconhecimento, considerar mais grave registrar o aluno errado (falso positivo)
do que deixar de registrar um aluno (falso negativo), já que este último ainda
é sinalizado ao motorista pelo retorno local de falha.

\textbf{Atividade 1.2 -- Revisão bibliográfica e avaliação empírica.} Além da
pesquisa bibliográfica, foi feito um benchmark com os dez modelos de
reconhecimento facial disponíveis na biblioteca DeepFace: VGG-Face
\cite{parkhi2015vggface}, FaceNet e FaceNet512 \cite{schroff2015facenet},
OpenFace \cite{amos2016openface}, DeepFace \cite{taigman2014deepface}, DeepID
\cite{sun2014deepid}, ArcFace \cite{deng2019arcface}, Dlib \cite{king2009dlib},
SFace \cite{zhong2021sface} e GhostFaceNet \cite{alansari2023ghostfacenets}.
O FaceNet512 é uma versão de 512 dimensões treinada pela comunidade, sem artigo
próprio; o modelo Dlib vem da biblioteca de mesmo nome.
A ferramenta foi escolhida porque aplica o mesmo pipeline de detecção e
alinhamento a todos os modelos, o que garante uma comparação justa. Cada
modelo foi avaliado em 384 pares de imagens da base LFW (96 pessoas, com pares
da mesma pessoa e de pessoas diferentes), medindo acurácia, precisão,
\textit{recall}, F1, tempo médio por par e tamanho do modelo. Cada modelo roda
em um processo separado, para que o tempo medido não seja distorcido pelos
modelos carregados antes. O benchmark foi rodado com dois detectores de rosto:
o do OpenCV, baseado no Haar Cascade \cite{viola2001haar}, mais leve, e o
RetinaFace \cite{deng2020retinaface}, mais preciso. Cada imagem
é detectada e alinhada uma única vez, e os dez modelos são avaliados sobre os
mesmos recortes, o que permite medir o tempo de detecção separado do tempo de
reconhecimento. A Tabela~\ref{tab:benchmark} mostra os modelos mais bem
colocados.

\begin{table}[H]
\centering
\caption{Resultados preliminares do benchmark (384 pares da base LFW, limiar padrão do DeepFace). O tempo por par inclui a detecção das duas imagens, medida na CPU de um computador pessoal.}
\label{tab:benchmark}
\begin{tabular}{lrrrrr}
    \toprule
     & \multicolumn{2}{c}{\textbf{OpenCV}} & \multicolumn{2}{c}{\textbf{RetinaFace}} & \\
    \cmidrule(lr){2-3} \cmidrule(lr){4-5}
    \textbf{Modelo} & \textbf{Acurácia} & \textbf{Tempo (s)} & \textbf{Acurácia} & \textbf{Tempo (s)} & \textbf{MB} \\
    \midrule
    ArcFace      & 92,97\% & 0,56 & 96,61\% & 1,80 & 130 \\
    VGG-Face     & 88,02\% & 1,06 & 95,57\% & 2,34 & 580 \\
    Dlib         & 84,64\% & 1,08 & 84,90\% & 2,30 & 22 \\
    Facenet      & 76,30\% & 0,64 & 82,55\% & 1,85 & 90 \\
    GhostFaceNet & 82,03\% & 0,51 & 81,77\% & 1,71 & 16 \\
    SFace        & 81,25\% & 0,09 & 79,69\% & 1,28 & 5 \\
    \bottomrule
\end{tabular}
\end{table}

O \textbf{ArcFace} é a escolha mais clara: tem a maior acurácia com os dois
detectores e chega a 96,61\% com o RetinaFace, acima da meta de 90\%, com
precisão de 99,45\%. O segundo modelo ainda está em aberto. O VGG-Face também
passa da meta com o RetinaFace (95,57\%), mas é o maior modelo avaliado
(580~MB). O SFace é de longe o mais leve e rápido (5~MB e 0,04~s de
reconhecimento por par), mas fica abaixo de 82\% com os dois detectores.

Os resultados também mostram três pontos que orientam os próximos passos.
Primeiro, o detector pesa tanto quanto o modelo de reconhecimento: o
RetinaFace leva cerca de 0,6~s por imagem na CPU, e com ele o ArcFace fica em
1,80~s por par, perto do limite de 2~s, o que torna a escolha do detector
parte da decisão. Segundo, SFace, GhostFaceNet e Dlib têm precisão próxima de
100\% e \textit{recall} baixo com qualquer detector, o que aponta para o
limiar padrão do DeepFace, conservador demais para esses modelos, e não para
uma limitação deles. Como o critério do projeto prioriza a precisão, a
comparação final vai fixar uma meta de precisão, calibrar o limiar de cada
modelo em um conjunto separado e comparar o \textit{recall} obtido nesse
ponto. Terceiro, com 384 pares a margem de erro é de cerca de 2~pontos
percentuais, de modo que ArcFace e VGG-Face ainda não se distinguem
estatisticamente. Até o fim de outubro, o benchmark será repetido com o
detector YuNet, os limiares serão calibrados, os dois modelos serão escolhidos
e a base de dados será ampliada para testá-los com mais rigor.

Uma limitação deve ser registrada: o benchmark mede verificação 1:1 entre
fotos de adultos, enquanto o sistema fará identificação entre cerca de 20
crianças, com a iluminação e o ângulo da porta do ônibus. Os números servem
para comparar os modelos entre si; o desempenho real só será conhecido nos
testes do TCC~II.

\textbf{Atividade 1.3 -- Projeto arquitetural e análise de trade-offs.} Nesta
atividade, será definida a arquitetura de hardware e software, com modelagem do
banco de dados e diagramas UML (casos de uso, componentes, sequência do fluxo de
embarque e implantação). Também será feita uma análise de trade-off entre
computação de borda com Raspberry Pi e microcontrolador ESP32, considerando
custo, capacidade de processamento e a possibilidade de rodar os modelos
escolhidos dentro do limite de 2~s. A tendência atual é usar Raspberry Pi~5
com o acelerador AI HAT (Hailo-10), já que o ESP32 não tem memória nem
capacidade de processamento para rodar com folga os modelos de reconhecimento
avaliados. Também serão especificados a câmera, o microfone direcional e o
banco de dados (PostgreSQL ou Firebase). Por fim, será estudado como emular o
ambiente do Hailo-10 no computador, para antecipar os testes de desempenho.
Os primeiros diagramas (contexto, casos de uso, componentes, implantação,
sequência do embarque e da sincronização e modelo de dados) já estão esboçados.
Para o software embarcado, foi adotada uma arquitetura de portas e adaptadores,
em que a câmera, o modelo de reconhecimento e o meio de envio ficam atrás de
interfaces e podem ser trocados sem alterar a regra de negócio, e uma fila
local de eventos com identificador único por evento, que garante o envio sem
perda nem duplicação quando a rede volta.

\textbf{Atividade 1.4 -- Avaliação de impactos.} Serão analisados os impactos
socioeconômicos (segurança física dos estudantes, transparência e otimização
dos recursos públicos), legais (tratamento de dados biométricos de menores de
idade segundo a LGPD, base legal, consentimento e prazos de retenção) e
ambientais (consumo de energia e descarte do hardware) da tecnologia proposta.

\textbf{Atividades 1.5 e 1.6 -- Monografia e seminário.} Redação do texto inicial
da monografia (Introdução, Revisão Bibliográfica, Metodologia, Avaliação de
Impactos e Cronograma do TCC~II), finalização do texto e preparação da
apresentação no Seminário de TCC~I.

\textbf{Ferramentas.} Python com OpenCV e TensorFlow Lite no pré-processamento
e inferência no dispositivo, biblioteca DeepFace na avaliação comparativa dos
modelos, MQTT na comunicação com o backend, PostgreSQL ou Firebase para
armazenamento e LaTeX/Overleaf para as documentações. Python e OpenCV foram
escolhidos pela maturidade em visão computacional e pelo suporte a hardware
embarcado. O TensorFlow Lite foi escolhido por ser otimizado para inferência
em dispositivos com poucos recursos. O MQTT, pelos motivos discutidos na
revisão bibliográfica. O Claude Code foi utilizado para estruturação da
documentação, acompanhamento do cronograma, geração de relatórios de status,
apoio no levantamento de questionamentos para os requisitos, apoio na revisão
do conteúdo produzido, tradução de algumas documentações e apoio no
entendimento, geração de resumos e indicação de fontes de pesquisa para serem
estudadas.

\subsection{Atividades do TCC II}

\begin{enumerate}
    \item \textbf{Implementação do software embarcado:} pipeline de captura,
          detecção, reconhecimento facial e de voz (incluindo a seleção e a
          avaliação do modelo de reconhecimento de locutor), retorno local
          por LED e bipe, e fila local de eventos.
    \item \textbf{Integração:} comunicação MQTT/API com o backend e o TRANSCOLAR
          RURAL, com sincronização offline sem perda nem duplicação, e módulo
          de cadastro e relatórios.
    \item \textbf{Simulação em ambiente controlado:} testes de acurácia e de
          tempo de resposta com o hardware definitivo, simulando as condições
          do ônibus (iluminação, movimento, perda de conexão).
    \item \textbf{Teste piloto:} instalação em um ônibus escolar rural real,
          com cerca de 20 alunos.
    \item \textbf{Validação e redação:} comparação das métricas obtidas com as
          metas (acurácia acima de 90\%, no máximo 2~s por aluno), redação da
          monografia de TCC~II e defesa.
\end{enumerate}

\subsection{Cronograma}

\begin{table}[H]
\centering
\caption{Cronograma do ciclo completo do TCC.}
\label{tab:cronograma}
\begin{tabularx}{\textwidth}{lX}
    \toprule
    \textbf{Data limite} & \textbf{Etapa} \\
    \midrule
    \multicolumn{2}{l}{\textit{TCC I}} \\
    08/2026     & 1.1 -- Levantamento de requisitos funcionais e não funcionais \\
    30/09/2026  & 1.2 -- Seleção de 10 algoritmos, implementação do benchmark e obtenção da base de dados de teste \\
    18/10/2026  & Entrega deste documento de Visão Geral \\
    31/10/2026  & 1.2 -- Benchmark com os dados iniciais, escolha dos 2 melhores algoritmos e ampliação da base de dados \\
    31/10/2026  & 1.3 -- Projeto arquitetural, modelagem de banco de dados e diagramas UML \\
    31/10/2026  & 1.4 -- Análise de impactos socioeconômicos, ambientais e de privacidade \\
    13/11/2026  & 1.5 -- Entrega do texto inicial da monografia \\
    16 a 19/11/2026 & 1.6 -- Seminários de TCC I \\
    30/11/2026  & 1.3 -- Estudo da emulação do Raspberry Pi AI HAT (Hailo-10) no computador \\
    \midrule
    \multicolumn{2}{l}{\textit{TCC II}} \\
    03/2027     & 2.1 -- Implementação do software embarcado \\
    04/2027     & 2.2 -- Integração MQTT/API com o TRANSCOLAR RURAL \\
    05/2027     & 2.3 -- Simulação em ambiente controlado \\
    06/2027     & 2.4 -- Teste piloto em ônibus escolar rural \\
    07/2027     & 2.5 -- Validação de métricas, monografia de TCC II e defesa \\
    \bottomrule
\end{tabularx}
\end{table}

% ============================================================================
% 4. REFERÊNCIAS
% ============================================================================
\bibliography{referencias}

\end{document}
